\documentclass[11pt,a4paper]{article}
\usepackage[utf8]{inputenc}
\usepackage[margin=1in]{geometry}
\usepackage{amsmath,amssymb,amsfonts}
\usepackage{graphicx}
\usepackage{booktabs}
\usepackage{hyperref}
\usepackage{float}
\usepackage{algorithm}
\usepackage{algpseudocode}

\hypersetup{
    colorlinks=true,
    linkcolor=black,
    citecolor=black,
    urlcolor=blue
}

\title{\textbf{Empirical Analysis of AdaHedge in Prediction with Expert Advice}}
\author{\textbf{Nitin Yadav}\\
M.Sc. in Data Science for Economics and Health\\
University of Milan\\
\texttt{nitin.yadav@studenti.unimi.it}}
\date{\today}

\begin{document}

\maketitle

\begin{abstract}
The Hedge algorithm is a standard method for prediction with expert advice, achieving a minimax regret bound of $\mathcal{O}(\sqrt{T \ln K})$. In practice, setting its learning rate $\eta$ requires knowing the horizon $T$ in advance. When the loss sequence is easy or stochastic, a fixed learning rate tuned for the worst-case makes the algorithm too slow to adapt, accumulating unnecessary regret. The AdaHedge algorithm (van Erven et al., 2011; de Rooij et al., 2014) addresses this by adjusting the learning rate $\eta_t = \frac{\ln K}{\Delta_{t-1}}$ using the cumulative mixability gap $\Delta_t$. In this project, we implement Hedge and AdaHedge from scratch in Python and compare their performance on three synthetic setups: stochastic, adversarial (shifting leader), and a low-gap case. Our simulation results show that AdaHedge achieves near-constant regret on stochastic data while remaining stable when the best expert changes over time.
\end{abstract}

\vspace{0.5em}
\noindent\textbf{Plagiarism Declaration:}
\textit{I declare that this material, which I now submit for assessment, is entirely my own work and has not been taken from the work of others, save and to the extent that such work has been cited and acknowledged within the text of my work. I understand that plagiarism, collusion, and copying are grave and serious offences in the university and accept the penalties that would be imposed should I engage in plagiarism, collusion or copying. This assignment, or any part of it, has not been previously submitted by me or any other person for assessment on this or any other course of study.}

\section{Introduction}

In online learning with expert advice, a learner makes sequential decisions over $T$ discrete rounds without assuming that the data comes from a stationary probability distribution.

At each round $t = 1, \dots, T$:
\begin{enumerate}
    \item The learner maintains a probability vector $\mathbf{w}_t = (w_{t,1}, \dots, w_{t,K})^\top$ over $K$ experts, with $w_{t,k} \ge 0$ and $\sum_{k=1}^K w_{t,k} = 1$.
    \item Nature reveals a loss vector $\boldsymbol{\ell}_t = (\ell_{t,1}, \dots, \ell_{t,K})^\top \in [0, 1]^K$.
    \item The learner incurs the expected loss $h_t = \mathbf{w}_t^\top \boldsymbol{\ell}_t = \sum_{k=1}^K w_{t,k} \ell_{t,k}$.
\end{enumerate}

The performance is evaluated through cumulative regret with respect to the best expert in hindsight:
\begin{equation}
    R_T = \sum_{t=1}^T h_t - \min_{k \in \{1, \dots, K\}} \sum_{t=1}^T \ell_{t,k}
\end{equation}
The goal of the learner is to guarantee sublinear regret, meaning $R_T / T \to 0$ as $T \to \infty$.

\section{Algorithms}

\subsection{Standard Hedge}
The standard Hedge algorithm (Freund and Schapire, 1997) updates expert weights exponentially based on cumulative losses:
\begin{equation}
    w_{t,k} = \frac{\exp(-\eta L_{t-1,k})}{\sum_{j=1}^K \exp(-\eta L_{t-1,j})}
\end{equation}
where $L_{t-1,k} = \sum_{s=1}^{t-1} \ell_{s,k}$ is the total loss of expert $k$ up to round $t-1$, and $\eta > 0$ is a fixed learning rate.

From standard theoretical analysis, setting $\eta^* = \sqrt{\frac{8 \ln K}{T}}$ yields an upper bound on regret:
\begin{equation}
    R_T \le \sqrt{\frac{T \ln K}{2}}
\end{equation}
This bound has two practical drawbacks:
\begin{enumerate}
    \item It requires knowing $T$ ahead of time to pick $\eta^*$.
    \item It assumes the worst-case. If one expert is consistently much better than the others (an easy stochastic setting), Follow-the-Leader would get $\mathcal{O}(1)$ regret, but Hedge with $\eta^*$ will still incur $\mathcal{O}(\sqrt{T})$ regret because its learning rate is forced to be small.
\end{enumerate}

\subsection{AdaHedge}
AdaHedge (van Erven et al., 2011; de Rooij et al., 2014) avoids tuning $\eta$ beforehand by tracking how difficult the sequence is round by round. 

For a given $\eta_t$, the mix loss $m_t$ is defined as:
\begin{equation}
    m_t(\eta_t) = -\frac{1}{\eta_t} \ln \left( \sum_{k=1}^K w_{t,k} \exp(-\eta_t \ell_{t,k}) \right)
\end{equation}
The difference between the expected loss $h_t$ and the mix loss $m_t$ is called the mixability gap:
\begin{equation}
    \delta_t = h_t - m_t(\eta_t)
\end{equation}
Because the logarithm is concave, Jensen's inequality guarantees that $\delta_t \ge 0$. In intuitive terms, $\delta_t$ measures how much the experts disagree. If the weights are concentrated on a single expert, or if all experts suffer the same loss, $\delta_t = 0$. When there is high variance in expert losses, $\delta_t > 0$.

AdaHedge accumulates these gaps over time, $\Delta_t = \sum_{s=1}^t \delta_s$, and sets the learning rate for the next round as:
\begin{equation}
    \eta_{t+1} = \frac{\ln K}{\Delta_t}
\end{equation}
At $t=1$, since $\Delta_0 = 0$, $\eta_1 = \infty$, which gives uniform weights initially.

\begin{algorithm}[H]
\caption{AdaHedge Implementation}
\begin{algorithmic}[1]
\State \textbf{Input:} Number of experts $K$, horizon $T$.
\State \textbf{Initialize:} $L_{0} = \mathbf{0} \in \mathbb{R}^K$, $\Delta_0 = 0$.
\For{$t = 1, \dots, T$}
    \State If $\Delta_{t-1} = 0$: set $\mathbf{w}_t = (1/K, \dots, 1/K)$ and $\eta_t = \infty$.
    \State Else: $\eta_t = \frac{\ln K}{\Delta_{t-1}}$, and $w_{t,k} \propto \exp(-\eta_t (L_{t-1,k} - \min_j L_{t-1,j}))$.
    \State Incur expected loss $h_t = \sum_{k=1}^K w_{t,k} \ell_{t,k}$.
    \State Observe $\boldsymbol{\ell}_t \in [0, 1]^K$.
    \State Compute mix loss $m_t = -\frac{1}{\eta_t} \ln \left( \sum_{k=1}^K w_{t,k} e^{-\eta_t \ell_{t,k}} \right)$.
    \State Compute gap $\delta_t = \max(0, h_t - m_t)$.
    \State Update $\Delta_t = \Delta_{t-1} + \delta_t$ and $L_{t,k} = L_{t-1,k} + \ell_{t,k}$.
\EndFor
\end{algorithmic}
\end{algorithm}

\section{Experimental Setup}

We ran our experiments with $K = 10$ experts over a horizon of $T = 5000$ rounds. Each test was repeated across 20 random seeds to average out noise.

\subsection{Datasets}
We tested the three synthetic environments specified in the assignment:
\begin{enumerate}
    \item \textbf{Stochastic:} Expert 0 has losses drawn from $\operatorname{Bernoulli}(0.3)$, while experts 1 through 9 have losses drawn from $\operatorname{Bernoulli}(0.5)$. Here, Expert 0 is clearly better ($\Delta p = 0.20$).
    \item \textbf{Adversarial (Shifting Leader):} The horizon $T$ is divided into 4 equal segments of 1250 rounds. In each segment, a different expert has $\operatorname{Bernoulli}(0.2)$ losses while the other nine have $\operatorname{Bernoulli}(0.8)$. This tests how quickly the algorithm shifts away from an expert that was previously good.
    \item \textbf{Low-Gap:} Expert 0 has losses from $\operatorname{Bernoulli}(0.49)$, and the others have losses from $\operatorname{Bernoulli}(0.50)$. The difference is subtle ($\Delta p = 0.01$), requiring thousands of samples to separate signal from noise.
\end{enumerate}

\subsection{Numerical Details}
In Python, calculating $\exp(-\eta L_k)$ can cause floating-point underflow or overflow when cumulative loss grows. We handled this by subtracting $\min_j L_j$ inside the exponential for both Hedge and AdaHedge. In the calculation of $m_t$, we also used the log-sum-exp shift trick and clipped exponential arguments to $[-700, 700]$.

\section{Results and Observations}

Table~\ref{tab:metrics} reports the average cumulative loss and final regret at $T = 5000$ across all 20 runs.

\begin{table}[H]
\centering
\caption{Average cumulative loss and final regret at $T = 5000$ (20 trials).}
\label{tab:metrics}
\begin{tabular}{llrrr}
\toprule
\textbf{Regime} & \textbf{Method} & \textbf{Cum. Loss} & \textbf{Best Expert} & \textbf{Final Regret $R_T$} \\
\midrule
\textbf{Stochastic} & Best Expert & 1503.20 & 1503.20 & 0.00 \\
& Hedge ($\eta = 0.005$) & 1953.33 & 1503.20 & 450.13 \\
& Hedge ($\eta^* = 0.061$) & 1543.06 & 1503.20 & 39.86 \\
& Hedge ($\eta = 0.500$) & 1510.38 & 1503.20 & 7.18 \\
& \textbf{AdaHedge} & \textbf{1509.16} & 1503.20 & \textbf{5.96} \\
\midrule
\textbf{Adversarial} & Best Expert & 3217.70 & 3217.70 & 0.00 \\
& Hedge ($\eta = 0.005$) & 3421.78 & 3217.70 & 204.08 \\
& Hedge ($\eta^* = 0.061$) & 3252.61 & 3217.70 & 34.91 \\
& Hedge ($\eta = 0.500$) & 3228.30 & 3217.70 & 10.60 \\
& \textbf{AdaHedge} & \textbf{3225.18} & 3217.70 & \textbf{7.48} \\
\midrule
\textbf{Low-Gap} & Best Expert & 2436.40 & 2436.40 & 0.00 \\
& Hedge ($\eta = 0.005$) & 2495.15 & 2436.40 & 58.75 \\
& Hedge ($\eta^* = 0.061$) & 2489.74 & 2436.40 & 53.34 \\
& Hedge ($\eta = 0.500$) & 2481.16 & 2436.40 & 44.76 \\
& \textbf{AdaHedge} & \textbf{2486.90} & 2436.40 & \textbf{50.50} \\
\bottomrule
\end{tabular}
\end{table}

\begin{figure}[H]
    \centering
    \includegraphics[width=\textwidth]{figures/fig1_regret_comparison.png}
    \caption{Cumulative Regret $R_t$ over time for the three settings (mean $\pm$ standard deviation over 20 runs).}
    \label{fig:regret}
\end{figure}

\begin{figure}[H]
    \centering
    \includegraphics[width=\textwidth]{figures/fig2_cum_loss_comparison.png}
    \caption{Cumulative Loss $L_t$ over time compared to the Best Single Expert in hindsight.}
    \label{fig:loss}
\end{figure}

\subsection{Discussion of Regret Curves}
\begin{itemize}
    \item \textbf{Stochastic case:} As shown in Figure~\ref{fig:regret}, AdaHedge quickly shifts its weight to Expert 0. Once this happens, the regret curve flattens out, giving a final regret of 5.96. In contrast, standard Hedge with minimax tuning $\eta^* = 0.061$ keeps accumulating regret because its learning rate is held down by the $\sqrt{T}$ schedule, finishing at $R_T = 39.86$. A very small learning rate ($\eta = 0.005$) performs poorly ($R_T = 450.13$) because it takes thousands of rounds to downweight bad experts.
    \item \textbf{Adversarial case:} When the best expert changes at rounds 1250, 2500, and 3750, the algorithm experiences sudden losses. AdaHedge adjusts its weights across each phase, ending with $R_T = 7.48$.
    \item \textbf{Low-Gap case:} Because the margin between 0.49 and 0.50 is small, variance dominates the first 1000 rounds. AdaHedge maintains stable performance without needing any manual tuning of $\eta$.
\end{itemize}

\begin{figure}[H]
    \centering
    \includegraphics[width=\textwidth]{figures/fig3_adahedge_dynamics.png}
    \caption{Evolution of the learning rate $\eta_t$ (left) and the cumulative mixability gap $\Delta_t$ (right).}
    \label{fig:dynamics}
\end{figure}

\begin{figure}[H]
    \centering
    \includegraphics[width=\textwidth]{figures/fig4_expert_weights.png}
    \caption{Weight evolution $w_{t,k}$ under AdaHedge in the stochastic setting (left) and adversarial setting (right).}
    \label{fig:weights}
\end{figure}

\subsection{Internal Dynamics of AdaHedge}
Figure~\ref{fig:dynamics} shows why AdaHedge works:
\begin{itemize}
    \item In the stochastic setting, once Expert 0 is identified, the weight on Expert 0 approaches 1 (Figure~\ref{fig:weights}, left). Because the algorithm predicts almost deterministically on the best expert, the mixability gap $\delta_t$ drops to zero. Consequently, $\Delta_t$ stops growing and flattens out around $\Delta_T \approx 5$. Since $\Delta_t$ stays constant, $\eta_t$ stays high ($\approx 0.4$), giving constant regret.
    \item In the adversarial setting, every time the leader changes, the algorithm makes wrong predictions, which causes $\delta_t > 0$. This pushes $\Delta_t$ up, which automatically lowers $\eta_t$. This acts as a brake, preventing the weights from oscillating erratically.
\end{itemize}

\section{Conclusion}
Our simulations confirm the practical benefits of AdaHedge. By tracking the cumulative mixability gap, the algorithm automatically adapts to the difficulty of the sequence: it gets constant $\mathcal{O}(1)$ regret on easy stochastic data, while preserving the worst-case $\mathcal{O}(\sqrt{T \ln K})$ guarantee on adversarial data, all without knowing the time horizon $T$ in advance.

\begin{thebibliography}{4}
\bibitem{vanerven2011}
T.~van Erven, P.~D. Gr{\"u}nwald, W.~M. Koolen, and S.~de~Rooij.
\newblock Adaptive Hedge.
\newblock In \emph{Advances in Neural Information Processing Systems (NeurIPS)}, volume~24, pages 1656--1664, 2011.

\bibitem{derooij2014}
S.~de~Rooij, T.~van Erven, P.~D. Gr{\"u}nwald, and W.~M. Koolen.
\newblock Follow the Leader If You Can, Hedge If You Must.
\newblock \emph{Journal of Machine Learning Research}, 15:1281--1316, 2014.

\bibitem{freund1997}
Y.~Freund and R.~E. Schapire.
\newblock A Decision-Theoretic Generalization of On-Line Learning and an Application to Boosting.
\newblock \emph{Journal of Computer and System Sciences}, 55(1):119--139, 1997.

\bibitem{cesabianchi2006}
N.~Cesa-Bianchi and G.~Lugosi.
\newblock \emph{Prediction, Learning, and Games}.
\newblock Cambridge University Press, 2006.
\end{thebibliography}

\end{document}
